\documentclass[10pt]{article}
\usepackage[a4paper,margin=0.88in]{geometry}
\renewcommand{\rmdefault}{cmr}
\usepackage{amsmath,amssymb,amsthm,mathtools,amscd,booktabs,enumitem,xcolor,hyperref}
\hypersetup{colorlinks=true,linkcolor=blue!55!black,citecolor=blue!55!black,urlcolor=blue!55!black}
\setlength{\emergencystretch}{3em}
\setlength{\parindent}{0pt}
\setlength{\parskip}{0.55em}
\setlist{leftmargin=1.6em,itemsep=0.25em,topsep=0.25em}
\sloppy

\newtheorem{theorem}{Theorem}[section]
\newtheorem{lemma}[theorem]{Lemma}
\newtheorem{proposition}[theorem]{Proposition}
\newtheorem{corollary}[theorem]{Corollary}
\theoremstyle{definition}
\newtheorem{definition}[theorem]{Definition}
\newtheorem{convention}[theorem]{Convention}
\newtheorem{remark}[theorem]{Remark}

\newcommand{\Fq}{\mathbf F_q}
\newcommand{\Zl}{\mathbf Z_\ell}
\newcommand{\Ql}{\mathbf Q_\ell}
\newcommand{\Qlb}{\overline{\mathbf Q}_\ell}
\newcommand{\OO}{\mathcal O}
\newcommand{\Aone}{\mathbf A^1}
\newcommand{\Db}{\mathcal D}
\newcommand{\Dbc}{\mathcal D^b_c}
\newcommand{\FT}{\operatorname{FT}}
\newcommand{\RHom}{R\mathcal Hom}
\newcommand{\Tr}{\operatorname{Tr}}
\newcommand{\Spec}{\operatorname{Spec}}
\newcommand{\Perv}{\operatorname{Perv}}

\begin{document}
\begin{center}
{\large\bfseries The Fourier--Deligne Transform\par}
\medskip
{\normalsize Artin--Schreier kernels, inversion, and perverse exactness\par}
\end{center}

\begin{abstract}
This note fixes the Fourier--Deligne transform with all maps, categories,
shifts, and Tate twists displayed.  It records its constructibility, base
change, inversion, duality, perversity, and convolution properties.  The
discussion takes place in ordinary triangulated derived categories.
\end{abstract}
\tableofcontents

\section{Coefficients and derived categories}

Fix a prime \(\ell\) and a finite extension \(E/\Ql\).  Let \(\OO_E\) be
the ring of integers of \(E\), choose a uniformizer \(\varpi\), and put
\(\OO_{E,n}=\OO_E/\varpi^n\OO_E\).  In this section \(X\) may be any
noetherian scheme of finite Krull dimension on which \(\ell\) is invertible.
The later six-functor statements concern separated morphisms of finite type
between schemes of this class.  No finite ground field is required until
geometric Frobenius is introduced.

\begin{definition}[Pro-etale coefficient sheaves]
For such a scheme \(X\), the small pro-etale site
\(X_{\mathrm{pro\acute et}}\) has weakly etale morphisms \(U\to X\) as
objects and the topology generated by jointly surjective fpqc families.
For a finite discrete ring \(\OO_{E,n}\), let \((\OO_{E,n})_X\) be the
sheafification of
\[
 U\longmapsto
 \operatorname{Cont}\!\left(\pi_0(|U|),\OO_{E,n}\right).
\]
Define sheaves of rings on \(X_{\mathrm{pro\acute et}}\) by
\[
 (\OO_E)_X=\varprojlim_n(\OO_{E,n})_X,
 \qquad
 E_X=(\OO_E)_X[1/\varpi]
     =\left(\varprojlim_n(\OO_{E,n})_X\right)[1/\varpi].
\]
Let \(\operatorname{Mod}((\OO_E)_X)\) and
\(\operatorname{Mod}(E_X)\) be the abelian categories of module sheaves
over these ring sheaves.  Their ordinary Verdier-localized derived
categories are
\[
 D(X,\OO_{E,n})=D\bigl(\operatorname{Mod}((\OO_{E,n})_X)\bigr),
 \qquad
 \Db(X,\OO_E)=D\bigl(\operatorname{Mod}((\OO_E)_X)\bigr),
 \qquad
 \Db(X,E)=D\bigl(\operatorname{Mod}(E_X)\bigr).
\]
Thus an object is represented by a cochain complex, and a morphism of
complexes is inverted precisely when it induces isomorphisms on all
cohomology sheaves.  Shifts and distinguished triangles are induced by
shifted complexes and mapping cones.
\end{definition}

\begin{definition}[Finite-level constructibility]
An \((\OO_{E,n})_X\)-module sheaf is constructible if \(X\) admits a
finite disjoint stratification by locally closed subschemes on each of which
the sheaf is locally constant with finite \(\OO_{E,n}\)-module geometric
stalks.  The category \(\Dbc(X,\OO_{E,n})\) is the full subcategory of
\(D(X,\OO_{E,n})\) consisting of complexes whose cohomology sheaves are
constructible and vanish outside a finite interval.
\end{definition}

\begin{definition}[Derived-complete integral category]
For \(K\in\Db(X,\OO_E)\), its \(n\)-th derived reduction is
\[
 K_n=K\otimes^L_{(\OO_E)_X}(\OO_{E,n})_X.
\]
The quotient maps form a tower, and derived inverse limit is obtained by
deriving inverse limit in the abelian category of towers of
\((\OO_E)_X\)-module sheaves.  The complex \(K\) is \emph{derived
\(\varpi\)-complete} if the canonical map
\[
 K\longrightarrow R\!\varprojlim_n K_n
\]
is an isomorphism in \(\Db(X,\OO_E)\).  Define
\(\Dbc(X,\OO_E)\) to be the full subcategory of derived
\(\varpi\)-complete complexes such that every
\(K_n\) lies in \(\Dbc(X,\OO_{E,n})\), one finite interval contains the
cohomological support of every \(K_n\), and, for \(m\le n\), the canonical
normalization map
\[
 K_n\otimes^L_{(\OO_{E,n})_X}(\OO_{E,m})_X
 \xrightarrow{\ \sim\ }K_m
\]
is an isomorphism.  This is a full ordinary triangulated subcategory of
\(\Db(X,\OO_E)\).
\end{definition}

\begin{definition}[Rationalization and integral lattice]
Derived extension of scalars along \((\OO_E)_X\to E_X\) is the exact
triangulated functor
\[
 \operatorname{rat}_E:\Dbc(X,\OO_E)\longrightarrow\Db(X,E),
 \qquad K\longmapsto K\otimes^L_{(\OO_E)_X}E_X.
\]
An \(\OO_E\)-lattice in \(L\in\Db(X,E)\) is an object
\(K\in\Dbc(X,\OO_E)\) together with an isomorphism
\[
 K\otimes^L_{(\OO_E)_X}E_X\xrightarrow{\ \sim\ }L.
\]
Let \({\Dbc}_{\mathrm{lat}}(X,E)\) be the full subcategory of objects that
admit such a lattice.
\end{definition}

\begin{definition}[Lisse, constructible, and finite Tor-amplitude]
An \(E_X\)-module sheaf is \emph{lisse of finite rank} if it is locally
isomorphic on the pro-etale site to \(E_X^d\) for some finite \(d\).  It is
\emph{constructible} if \(X\) has a finite stratification by locally closed
constructible subschemes on each of which its restriction is lisse of finite
rank.  Let \({\Dbc}_{\mathrm{int}}(X,E)\) be the full subcategory of
\(\Db(X,E)\) consisting of complexes \(L\) for which only finitely many
\(\mathcal H^i(L)\) are nonzero and every \(\mathcal H^i(L)\) is
constructible.  This is the intrinsic bounded-constructible condition.

A complex \(L\) has finite Tor-amplitude if there are integers \(a\le b\)
such that, for every \(E_X\)-module sheaf \(M\) in degree zero,
\(\mathcal H^i(L\otimes_{E_X}^LM)=0\) for \(i\notin[a,b]\).
\end{definition}

\begin{theorem}[Integral-lattice and intrinsic comparison]
The subcategories \({\Dbc}_{\mathrm{lat}}(X,E)\) and
\({\Dbc}_{\mathrm{int}}(X,E)\) coincide.  Their common category is denoted
\(\Dbc(X,E)\).  It is an idempotent-complete full triangulated subcategory
of \(\Db(X,E)\), every one of its objects has finite Tor-amplitude, and
every object admits an \(\OO_E\)-lattice.

More precisely, let \(\Dbc(X,\OO_E)[1/\varpi]\) have the same objects as
\(\Dbc(X,\OO_E)\) and morphism groups
\[
 \operatorname{Hom}(K,L)\otimes_{\OO_E}E.
\]
Its distinguished triangles are those isomorphic to rationalizations of
distinguished triangles, and \((-)^\natural\) denotes its Karoubi
completion.  Rationalization induces an exact equivalence
\[
 \Dbc(X,\OO_E)[1/\varpi]^\natural
 \xrightarrow{\ \sim\ }\Dbc(X,E).
\]
Under this equivalence, reduction of an integral lattice modulo
\(\varpi^n\) is its derived tensor product with
\((\OO_{E,n})_X\).  In particular the definition is independent, up to
canonical equivalence, of the chosen uniformizer.
\end{theorem}

\begin{definition}[The \(\Qlb\)-coefficient category]
Let finite extensions \(E/\Ql\) be taken inside \(\Qlb\).  First form the
finite-stage scalar-extension category \(\mathcal C_X\).  Its objects are
pairs \((E,K)\) with \(K\in\Dbc(X,E)\), and its morphisms are
\[
 \operatorname{Hom}_{\mathcal C_X}((E,K),(F,L))
 =\varinjlim_{M\supset E,F}
 \operatorname{Hom}_{\Dbc(X,M)}
 \left(K\otimes_E^LM,L\otimes_F^LM\right),
\]
where \(M\) runs over finite extensions of \(\Ql\) in \(\Qlb\) containing
both coefficient fields.  Composition is performed after passing all
representatives to one common finite extension; associativity follows from
the associativity of finite scalar extension.  The shift of \((E,K)\) is
\((E,K[1])\), and a triangle in \(\mathcal C_X\) is distinguished if it is
isomorphic to the image of a distinguished triangle in \(\Dbc(X,M)\) for
some common finite extension \(M/\Ql\).  These operations make
\(\mathcal C_X\) an ordinary triangulated category.

Define \(\Dbc(X,\Qlb)=\mathcal C_X^\natural\), the Karoubi completion of
this ordinary triangulated category.  Concretely, an object is a triple
\((E,K,e)\), where \((E,K)\in\mathcal C_X\) and
\(e\in\operatorname{End}_{\mathcal C_X}(E,K)\) satisfies \(e^2=e\).  A
morphism
\[
 (E,K,e)\longrightarrow(F,L,e')
\]
is a class \(u\) in the preceding common-finite-stage colimit satisfying
\(u=e'u=ue\).  Shifts are \((E,K[1],e[1])\), and the distinguished
triangles are the triangles isomorphic to direct summands of distinguished
triangles represented over a common finite extension of \(\Ql\).  This makes
\(\Dbc(X,\Qlb)\) an ordinary idempotent-complete triangulated category.
\end{definition}

\begin{theorem}[Operations after scalar extension]
For every finite extension \(E'/E\), derived scalar extension
\[
 -\otimes_E^LE':\Dbc(X,E)\longrightarrow\Dbc(X,E')
\]
is exact and is compatible with composition of coefficient extensions.
Every derived operation used in this note commutes with finite scalar
extension through its canonical comparison isomorphism.  It therefore acts
on \(\mathcal C_X\), preserves idempotents, and extends uniquely to the
Karoubi triples defining \(\Dbc(X,\Qlb)\).
\end{theorem}

In formulas, \(\otimes_E^L\) abbreviates tensor product over the coefficient
sheaf \(E_X\) on the scheme in question, and the symbol \(E\) used as a sheaf means
\(E_X\).

\begin{definition}[Ambient derived operations]
\label{def:ambient-operations}
For a morphism \(f:X\to Y\), inverse image and direct image of sheaves induce
\[
 Lf^*:\Db(Y,E)\longrightarrow\Db(X,E),
 \qquad
 Rf_*:\Db(X,E)\longrightarrow\Db(Y,E).
\]
Here \(Lf^*\) is computed by a \(K\)-flat resolution and \(Rf_*\) by a
\(K\)-injective resolution.  The derived tensor product
\[
 -\otimes_E^L-:\Db(X,E)\times\Db(X,E)\longrightarrow\Db(X,E)
\]
is likewise computed using a \(K\)-flat resolution.
The derived internal Hom is the bifunctor
\[
 R\mathcal Hom_X(-,-):
 \Db(X,E)^{\mathrm{op}}\times\Db(X,E)\longrightarrow\Db(X,E),
\]
computed using a \(K\)-injective resolution in the second variable.  It
restricts on bounded constructible complexes to
\[
 R\mathcal Hom_X(-,-):
 \Dbc(X,E)^{\mathrm{op}}\times\Dbc(X,E)\longrightarrow\Dbc(X,E).
\]
For the structural morphism \(a_X:X\to\Spec k\) of a finite-type
\(k\)-scheme, the global derived Hom complex is
\[
 R\operatorname{Hom}_X(K,L)
 =R\Gamma\!\left(X,R\mathcal Hom_X(K,L)\right)
 \quad\text{in }D(E),
\]
where \(D(E)\) is the ordinary derived category of \(E\)-vector spaces.

If \(f\) is separated and of finite type, choose a compactification
\(X\xrightarrow{j}\overline X\xrightarrow{\bar f}Y\), with \(j\) open and
\(\bar f\) proper, and define
\[
 Rf_!=R\bar f_*\,j_!:\Db(X,E)\longrightarrow\Db(Y,E),
\]
where \(j_!\) is extension by zero.  The compactification-independence
theorem canonically identifies the results for two choices, compatibly with
identities and composition.  Restrictions of these ambient functors to
bounded constructible categories are always justified by the finiteness
theorems stated below.
\end{definition}

\begin{definition}[Cohomological shifts and Tate twists]
For an integer \(n\), \(K[n]\) denotes cohomological shift.  On
\(X_{\mathrm{pro\acute et}}\), define
\[
 (\Zl(1))_X=\varprojlim_m\mu_{\ell^m,X},
\]
where the transition map
\(\mu_{\ell^{m+1},X}\to\mu_{\ell^m,X}\) is \(z\mapsto z^\ell\).  Put
\[
 (\OO_E(1))_X=(\Zl(1))_X
   \otimes_{(\Zl)_X}(\OO_E)_X,
 \qquad
 E(1)_X=(\OO_E(1))_X\otimes_{(\OO_E)_X}E_X.
\]
For \(n\geq0\), let \(E(n)_X=E(1)_X^{\otimes_E n}\); for \(n<0\), let
\(E(n)_X=(E(-n)_X)^\vee\), where
\((-)^\vee=\mathcal Hom_{E_X}(-,E_X)\).  The \(n\)-th Tate twist is
\[
 K(n)=K\otimes_E^LE(n)_X.
\]
It is an exact triangulated functor.  The same construction at the integral
and finite coefficient levels is compatible with derived reduction,
rationalization, finite scalar extension, and the \(\Qlb\)-coefficient
category.
\end{definition}

Fix from now on a prime power \(q=p^a\) with \(\ell\ne p\), put
\(k_0=\Fq\) and \(\bar k=\overline{\Fq}\), and use a subscript \(0\) for
schemes over \(k_0\).  For \(X_0/k_0\), its geometric base change is always
written
\[
 X=X_0\times_{k_0}\bar k.
\]

\begin{definition}[Geometric Frobenius]
Let
\[
 \sigma_{\mathrm{ar}}:\overline{\Fq}\longrightarrow\overline{\Fq},
 \quad a\longmapsto a^q,
 \qquad
 \sigma_{\mathrm{geom}}=\sigma_{\mathrm{ar}}^{-1}.
\]
For \(X_0/k_0\) and its geometric base change \(X\), let
\(\gamma_X:X\to X\) denote the semilinear automorphism induced by
\(\sigma_{\mathrm{geom}}\).  The \(q\)-power scheme Frobenius
\(\operatorname{Fr}_{X_0,q}:X_0\to X_0\), locally given by
\(f\mapsto f^q\), is a different morphism and is not denoted by
\(\gamma_X\).
\end{definition}

\begin{definition}[Strict equivariant complexes]
Let \(\mathsf C^b_{W,c}(X_0,E)\) be the category whose objects are pairs
\((C^\bullet,\Phi_C)\) consisting of a bounded cochain complex of
\(E_X\)-module sheaves with constructible cohomology and a chain
isomorphism
\[
 \Phi_C:\gamma_X^*C^\bullet\xrightarrow{\ \sim\ }C^\bullet.
\]
A morphism
\(u:(C^\bullet,\Phi_C)\to(D^\bullet,\Phi_D)\) is a chain map satisfying
\[
 u\circ\Phi_C=\Phi_D\circ\gamma_X^*u.
\]
Two such maps are homotopic if they are joined by a chain homotopy satisfying
the same equivariance identity.  The resulting homotopy category is denoted
\(\mathsf K^b_{W,c}(X_0,E)\).  Shifts use the shifted equivariant chain
isomorphism, and the cone of an equivariant chain map carries the direct-sum
equivariant chain isomorphism.  Hence \(\mathsf K^b_{W,c}(X_0,E)\) is an
ordinary triangulated category.

An equivariant chain map is an \emph{equivariant quasi-isomorphism} if its
underlying chain map induces an isomorphism on every cohomology sheaf.  The
bounded constructible Weil derived category is the Verdier localization
\[
 \mathcal D^b_{W,c}(X_0,E)
 =\mathsf K^b_{W,c}(X_0,E)
  [\{\text{equivariant quasi-isomorphisms}\}^{-1}].
\]
This is the meaning of a Weil complex in this note.  In particular it is not
defined by first choosing an object of \(\Dbc(X,E)\) and then adjoining an
isomorphism inside that already localized category.
\end{definition}

\begin{definition}[Strict Weil complexes with \(\Qlb\)-coefficients]
For a finite extension \(E'/E\) inside \(\Qlb\), tensoring the underlying
strict complex and its equivariant chain isomorphism with \(E'\) defines an
exact scalar-extension functor
\[
 -\otimes_E^LE':\mathcal D^b_{W,c}(X_0,E)
 \longrightarrow\mathcal D^b_{W,c}(X_0,E').
\]
Let \(\mathcal C^W_{X_0}\) be the finite-stage category whose objects are
pairs \((E,\mathcal K)\) with
\(\mathcal K\in\mathcal D^b_{W,c}(X_0,E)\), and whose morphisms are
\[
 \operatorname{Hom}_{\mathcal C^W_{X_0}}
 ((E,\mathcal K),(F,\mathcal L))
 =\varinjlim_{M\supset E,F}
 \operatorname{Hom}_{\mathcal D^b_{W,c}(X_0,M)}
 (\mathcal K\otimes_E^LM,\mathcal L\otimes_F^LM).
\]
Here \(M\) runs over finite extensions of \(\Ql\) in \(\Qlb\) containing
both coefficient fields.  Composition is performed over a common finite
stage.  Shifts and distinguished triangles are represented over a common
finite stage, making \(\mathcal C^W_{X_0}\) an ordinary triangulated
category.

Define
\[
 \mathcal D^b_{W,c}(X_0,\Qlb)=(\mathcal C^W_{X_0})^\natural.
\]
Thus an object is a Karoubi triple \((E,\mathcal K,e)\), with
\(e^2=e\) in the common-stage endomorphism colimit.  A morphism from
\((E,\mathcal K,e)\) to \((F,\mathcal L,e')\) is a common-stage class
\(u\) satisfying \(u=e'u=ue\).  Shifts act on all three entries, and
distinguished triangles are the direct summands of finite-stage
distinguished triangles.  This is an ordinary idempotent-complete
triangulated category.
\end{definition}

\begin{definition}[Forgetful functor and Frobenius operator]
Forgetting the strict equivariant chain isomorphism preserves cones and
quasi-isomorphisms, and therefore induces an exact triangulated functor
\[
 \operatorname{For}_X:\mathcal D^b_{W,c}(X_0,E)
 \longrightarrow\Dbc(X,E).
\]
For \(\mathcal K\in\mathcal D^b_{W,c}(X_0,E)\), write
\(K=\operatorname{For}_X(\mathcal K)\).  Its strict datum induces an
isomorphism \(\varphi:\gamma_X^*K\xrightarrow{\sim}K\) in
\(\Dbc(X,E)\), but the notation \((K,\varphi)\) records only the image of
the strictly equivariant object under \(\operatorname{For}_X\).

For \(n\geq1\), define the iterated chain isomorphism
\[
 \Phi_C^{(n)}=\Phi_C\circ\gamma_X^*\Phi_C\circ\cdots\circ
 (\gamma_X^{n-1})^*\Phi_C:
 (\gamma_X^n)^*C^\bullet\xrightarrow{\ \sim\ }C^\bullet.
\]
If \(x\in X_0(\mathbf F_{q^n})\), choose the associated geometric point
\(\bar x\to X\).  Since \(\gamma_X^n\bar x=\bar x\), the stalk of
\(\Phi_C^{(n)}\) defines
\[
 F_{x,n}:\mathcal H^i(K)_{\bar x}\longrightarrow
          \mathcal H^i(K)_{\bar x}.
\]
Its conjugacy class is independent of the chosen representative of the
geometric point.  If \(x\) is a closed point of degree \(d\), write
\(F_x=F_{x,d}\); if \(n=md\), then \(F_{x,n}=F_x^m\).  Every operator in
this convention is geometric Frobenius.

Finite-stage forgetful functors and scalar extension induce an exact
triangulated functor
\[
 \operatorname{For}_X:\mathcal D^b_{W,c}(X_0,\Qlb)
 \longrightarrow\Dbc(X,\Qlb),
\]
which sends \((E,\mathcal K,e)\) to
\((E,\operatorname{For}_X\mathcal K,\operatorname{For}_X(e))\).
A bounded
constructible complex on \(X_0\) pulls back to a strict Weil complex on
\(X\) through its canonical descent isomorphism.
\end{definition}

\begin{proposition}[Tate twists as strict Weil complexes]
The sheaf \(E(1)_X\) is the geometric pullback of the arithmetic Tate
sheaf and carries its canonical strict equivariant chain isomorphism.  For
\(x\in X_0(\mathbf F_{q^n})\), the operator \(F_{x,n}\) on its stalk is
multiplication by \(q^{-n}\).  Tensoring with this object defines Tate
twists on \(\mathcal D^b_{W,c}(X_0,E)\); finite scalar extension and
Karoubi completion define the same operation on
\(\mathcal D^b_{W,c}(X_0,\Qlb)\).
\end{proposition}

\begin{theorem}[Constructibility of the direct-image operations]
\label{thm:constructible-direct-images}
Let \(k\) be either \(k_0\) or \(\bar k\), and let \(f:X\to Y\) be a
separated morphism of finite type between separated schemes of finite type
over \(k\).  The ambient operations of
Definition~\ref{def:ambient-operations} restrict to exact triangulated
functors
\[
 \begin{aligned}
 Lf^*&:\Dbc(Y,E)\longrightarrow\Dbc(X,E),\\
 Rf_*&:\Dbc(X,E)\longrightarrow\Dbc(Y,E),\\
 Rf_!&:\Dbc(X,E)\longrightarrow\Dbc(Y,E).
 \end{aligned}
\]
The restrictions commute with finite extension of the coefficient field.
\end{theorem}

\begin{definition}[Exceptional-adjoint datum]
\label{def:exceptional-adjoint}
Let \(f:X\to Y\) be a separated morphism of finite type between the schemes
considered above.  An \emph{exceptional-adjoint datum} for \(f\) consists of
an exact triangulated functor
\[
 Rf^!:\Dbc(Y,E)\longrightarrow\Dbc(X,E)
\]
and natural transformations
\[
 \eta_f:\operatorname{id}\longrightarrow Rf^!Rf_!,
 \qquad
 \varepsilon_f:Rf_!Rf^!\longrightarrow\operatorname{id}
\]
satisfying the two triangular identities.  Equivalently, \(Rf^!\) is a
right adjoint of the restriction
\[
 Rf_!:\Dbc(X,E)\longrightarrow\Dbc(Y,E)
\]
of the compactly supported direct image defined above.  The datum includes
the displayed source and target for \(Rf^!\); constructibility of \(Rf_!\)
is Theorem~\ref{thm:constructible-direct-images}.  No right adjoint on the
entire unbounded ambient category is part of this definition.
\end{definition}

\begin{definition}[Verdier-admissible schemes and duality]
\label{def:verdier-admissible}
Let \(k\) be either \(k_0\) or \(\bar k\), let \(X\) be separated and of
finite type over \(k\), and write \(a_X:X\to\Spec k\) for its structural
morphism.  A \emph{Verdier-admissibility datum} for \(X\) is an
exceptional-adjoint datum for \(a_X\), together with the assertions
\[
 \omega_X:=Ra_X^!E_k\in\Dbc(X,E)
\]
where \(E_k\) denotes the constant coefficient sheaf on \(\Spec k\),
and the assertion that \(R\mathcal Hom_X(-,\omega_X)\) preserves
\(\Dbc(X,E)\).  For such a datum, define Verdier duality by
\[
 D_X(K):=R\mathcal Hom_X(K,\omega_X):
 \Dbc(X,E)^{\mathrm{op}}\longrightarrow\Dbc(X,E).
\]
This is an internal-Hom sheaf complex on \(X\), not the global complex
\(R\operatorname{Hom}_X(K,\omega_X)\in D(E)\).
\end{definition}

\begin{theorem}[Equivariant six operations]
Let \(f_0:X_0\to Y_0\) be a \(k_0\)-morphism and let
\(f:X\to Y\) be its geometric base change.  Assume the bounded-
constructibility conditions for the displayed operations, fix an
exceptional-adjoint datum for \(f\), and fix Verdier-admissibility data for
\(X\) and \(Y\).  There are exact triangulated functors on Weil derived
categories
\[
 \begin{aligned}
 Lf^*&:\mathcal D^b_{W,c}(Y_0,E)
       \longrightarrow\mathcal D^b_{W,c}(X_0,E),\\
 Rf_*,\ Rf_!&:\mathcal D^b_{W,c}(X_0,E)
       \longrightarrow\mathcal D^b_{W,c}(Y_0,E),\\
 Rf^!&:\mathcal D^b_{W,c}(Y_0,E)
       \longrightarrow\mathcal D^b_{W,c}(X_0,E).
 \end{aligned}
\]
Here \(Rf_!\) is used for separated finite-type \(f_0\), and \(Rf^!\) is the
right adjoint specified by Definition~\ref{def:exceptional-adjoint}.
Derived tensor product, derived internal Hom, Tate twists, and the already
defined dualities \(D_X,D_Y\) have corresponding exact lifts.  The
forgetful functors intertwine
each lift with the corresponding operation on \(\Dbc(X,E)\), through a
canonical natural isomorphism compatible with identities, composition,
adjunction, tensor--Hom adjunction, and base change.
For \(\Qlb\)-coefficients these functors are defined on finite-stage pairs,
their morphisms are the induced maps in the common-stage Hom colimits, and
their values on Karoubi triples are obtained by applying the functor to the
idempotent.  They give the same six operations, tensor product, internal
Hom, Tate twists, and Verdier duality on
\(\mathcal D^b_{W,c}(-,\Qlb)\), compatibly with the \(\Qlb\)-forgetful
functors.
\end{theorem}

\section{The Artin--Schreier sheaf}

\begin{definition}[Additive character datum]
Fix a nontrivial homomorphism
\[
  \psi:(\Fq,+)\longrightarrow E^\times
\]
and assume that \(E\) contains all values of \(\psi\).  Its inverse
character is
\(\psi^{-1}(a)=\psi(-a)=\psi(a)^{-1}\).  Let
\(\widehat{\Fq}=\operatorname{Hom}((\Fq,+),E^\times)\).
\end{definition}

\begin{definition}[Artin--Schreier cover and translations]
Write \(X_0=\Spec k_0[z]\), \(Y_0=\Spec k_0[t]\), and define
\[
 \operatorname{AS}_q:X_0\longrightarrow Y_0,
 \qquad t\longmapsto z^q-z.
\]
For \(a\in\Fq\), let \(\tau_a:X_0\to X_0\) be given by
\(z\mapsto z+a\).  The equality
\((z+a)^q-(z+a)=z^q-z\) makes every \(\tau_a\) an automorphism over
\(Y_0\), and \(\tau_a\tau_b=\tau_{a+b}\).
\end{definition}

\begin{lemma}[Coordinate presentation and finite freeness]
The coordinate algebra of \(X_0\) has the presentation
\[
 k_0[z]\simeq k_0[t][Z]/(Z^q-Z-t).
\]
It is a finite free \(k_0[t]\)-module of rank \(q\), with basis
\(1,z,\ldots,z^{q-1}\).  Hence \(\operatorname{AS}_q\) is finite locally
free of constant degree \(q\).
\end{lemma}

\begin{theorem}[The finite-etale Artin--Schreier torsor]
The morphism \(\operatorname{AS}_q\) is etale.  The translation action of
the constant group scheme \((\Fq)_{Y_0}\) is free, and the canonical map
\[
 (\Fq)_{Y_0}\times_{Y_0}X_0
 \longrightarrow X_0\times_{Y_0}X_0,
 \qquad (a,z)\longmapsto(z,z+a),
\]
is an isomorphism.  The invariant subring is
\(k_0[z]^{\Fq}=k_0[z^q-z]\), so \(Y_0\) is the affine categorical quotient
of \(X_0\) by \(\Fq\).  These assertions are compatible with arithmetic
and geometric Frobenius.
\end{theorem}

\begin{definition}[Translation operators and character projectors]
Put
\[
 P=(\operatorname{AS}_q)_*E_{X_0},
\]
a lisse sheaf of rank \(q\) on \(Y_0\).  Translation by \(a\) induces
\(T_a:P\to P\), with \(T_0=1\) and \(T_aT_b=T_{a+b}\).  Since \(q\) is
invertible in \(E\), define, for \(\chi\in\widehat{\Fq}\),
\[
 e_\chi=\frac1q\sum_{a\in\Fq}\chi(a)^{-1}T_a
       \in\operatorname{End}(P).
\]
\end{definition}

\begin{lemma}[Character-projector algebra]
For \(\chi,\chi'\in\widehat{\Fq}\),
\[
 e_\chi e_{\chi'}=\begin{cases}e_\chi,&\chi=\chi',\\0,&\chi\ne\chi',
 \end{cases}
 \qquad
 \sum_{\chi\in\widehat{\Fq}}e_\chi=1_P,
 \qquad
 T_b e_\chi=\chi(b)e_\chi.
\]
Consequently the image and kernel of every \(e_\chi\) are lisse direct
summands of \(P\), and
\(P\simeq\bigoplus_{\chi\in\widehat{\Fq}}\operatorname{im}(e_\chi)\).
\end{lemma}

\begin{definition}[Artin--Schreier character sheaf]
The arithmetic Artin--Schreier sheaf attached to \(\psi\) is
\[
 \mathcal L_{\psi,0}=\operatorname{im}(e_\psi)\subset P.
\]
The inclusion \(s_\psi:\mathcal L_{\psi,0}\to P\) and retraction
\(r_\psi:P\to\mathcal L_{\psi,0}\) are fixed by
\[
 r_\psi s_\psi=1,\qquad s_\psi r_\psi=e_\psi.
\]
Let \(\mathcal L_\psi\) be the pullback to
\(\Aone=Y_0\times_{k_0}\bar k\).  Frobenius commutes with every
\(T_a\), hence preserves \(e_\psi\) and gives
\(\mathcal L_\psi\) its strict Weil structure.
\end{definition}

\begin{theorem}[Rank, lissity, and Frobenius trace]
The sheaf \(\mathcal L_{\psi,0}\) is lisse of rank one.  For every
\(n\ge1\) and every \(x\in\mathbf F_{q^n}\), geometric Frobenius on the
stalk at \(x\) acts with trace
\[
 \Tr\!\left(F_{x,n}\mid(\mathcal L_{\psi,0})_{\bar x}\right)
 =\psi\!\left(\Tr_{\mathbf F_{q^n}/\Fq}(x)\right).
\]
Thus \(\mathcal L_{\psi,0}\) is punctually pure of weight zero for every
embedding of its coefficient field into \(\mathbf C\).
\end{theorem}

\begin{definition}[Rigidification at the origin]
The point \(z=0\) above \(t=0\), together with the character vector in the
regular representation of \(\Fq\), gives a Frobenius-equivariant
isomorphism
\[
 \epsilon_\psi:E\xrightarrow{\sim}0^*\mathcal L_{\psi,0}.
\]
This is the chosen rigidification of \(\mathcal L_{\psi,0}\).
\end{definition}

\begin{theorem}[Multiplicative structure for addition]
Let \(m:\Aone_0\times\Aone_0\to\Aone_0\) be addition and let
\(u_1,u_2\) be the projections.  There is a unique
Frobenius-equivariant isomorphism
\[
 c_\psi:m^*\mathcal L_{\psi,0}
 \xrightarrow{\sim}
 u_1^*\mathcal L_{\psi,0}\otimes_Eu_2^*\mathcal L_{\psi,0}
\]
whose restriction along \((0,0)\) agrees with \(\epsilon_\psi\).  The
isomorphism \(c_\psi\) satisfies the left and right unit identities,
symmetry under interchange of the two factors, and the two associativity
identities on \((\Aone_0)^3\).  Its geometric pullback has the same
properties.
\end{theorem}

\begin{theorem}[Inverse-character duality]
Evaluation through the regular-representation pairing restricts to a
perfect Frobenius-equivariant pairing
\[
 \mathcal L_{\psi,0}\otimes_E\mathcal L_{\psi^{-1},0}
 \longrightarrow E_{\Aone_0}.
\]
It induces canonical isomorphisms
\[
 \mathcal L_{\psi,0}^{\vee}\simeq\mathcal L_{\psi^{-1},0},
 \qquad
 \mathcal L_\psi^{\vee}\simeq\mathcal L_{\psi^{-1}}.
\]
They are compatible with the rigidifications and with \(c_\psi\).
\end{theorem}

\begin{proposition}[Extension of the finite ground field]
For \(m\ge1\), put
\[
 \psi_m=\psi\circ\Tr_{\mathbf F_{q^m}/\Fq}.
\]
For the projection \(b_m:\Aone_{\mathbf F_{q^m}}\to\Aone_0\), there is a
canonical Frobenius-equivariant isomorphism
\[
 b_m^*\mathcal L_{\psi,0}\xrightarrow{\sim}\mathcal L_{\psi_m,0}.
\]
It respects rigidification, addition, inverse-character duality, and every
finite extension of the coefficient field.  Consequently finite ground
field extension carries the Fourier transform for \(\psi\) to the Fourier
transform for \(\psi_m\).
\end{proposition}

\section{Vector bundles and the two transforms}

Let the ground field \(k\) be either \(k_0\) or \(\bar k\), let \(S\) be a
separated finite-type \(k\)-scheme, and let \(\mathcal V\) be a locally free
\(\mathcal O_S\)-module of constant rank \(r\).  The notation
\(\mathcal L_{\psi,S}\) has the arithmetic or geometric meaning specified
above according to \(k\).  Define the total spaces
\[
 V=\Spec_S\operatorname{Sym}_{\mathcal O_S}(\mathcal V^\vee),
 \qquad
 V^\vee=\Spec_S\operatorname{Sym}_{\mathcal O_S}(\mathcal V).
\]
Thus \(V^\vee\) is the dual vector bundle.  Evaluation of a vector by a
covector gives an \(S\)-morphism
\[
 \mu=\langle\ ,\ \rangle:V\times_SV^\vee\longrightarrow\Aone_S.
\]
Write
\[
 p_1:V\times_SV^\vee\longrightarrow V,
 \qquad
 p_2:V\times_SV^\vee\longrightarrow V^\vee
\]
for the projections and put
\(\mathcal K_\psi=L\mu^*\mathcal L_{\psi,S}\).

\begin{definition}[Fourier integrand]
For \(K\in\Db(V,E)\), its Fourier integrand is
\[
 \mathcal A_\psi(K)
 =Lp_1^*K\otimes_E^L\mathcal K_\psi
 \quad\text{in }\Db(V\times_SV^\vee,E).
\]
\end{definition}

\begin{definition}[Fourier--Deligne transforms]
The compactly supported and ordinary Fourier--Deligne transforms are the
triangulated functors
\[
 \begin{aligned}
 \FT^{V/S}_{\psi,!}(K)
   &=Rp_{2!}\mathcal A_\psi(K)[r],\\
 \FT^{V/S}_{\psi,*}(K)
   &=Rp_{2*}\mathcal A_\psi(K)[r].
 \end{aligned}
\]
Both have ambient source \(\Db(V,E)\) and target \(\Db(V^\vee,E)\).  The
shift \([r]\) belongs to the definition; there is no Tate twist in the
definition.  Unless a star is displayed, \(\FT_\psi\) means
\(\FT_{\psi,!}\).
\end{definition}

The canonical forget-supports morphism \(Rp_{2!}\to Rp_{2*}\) gives a
natural transformation
\[
 \alpha_K:\FT_{\psi,!}(K)\longrightarrow\FT_{\psi,*}(K).
\]

\begin{theorem}[Weil lift of the Fourier--Deligne transform]
Let \(S_0\) be a separated finite-type \(k_0\)-scheme, let
\(V_0\to S_0\) be a vector bundle, and write \(S\), \(V\), and \(V^\vee\)
for the geometric base changes of \(S_0\), \(V_0\), and \(V_0^\vee\).
The canonical Weil datum on
\(\mathcal L_\psi\) pulls back along \(\mu\) to a strict Weil datum on
\(\mathcal K_\psi\).  The lifted pullback, tensor product, and direct-image
functors of the preceding theorem therefore define exact triangulated
functors
\[
 \begin{aligned}
 \FT^{V_0/S_0}_{\psi,!},\ \FT^{V_0/S_0}_{\psi,*}:
 \mathcal D^b_{W,c}(V_0,E)&\longrightarrow
 \mathcal D^b_{W,c}(V_0^\vee,E),\\
 \mathcal K&\longmapsto
 Rp_{2!}\bigl(Lp_1^*\mathcal K\otimes_E^L\mathcal K_\psi\bigr)[r],\\
 \mathcal K&\longmapsto
 Rp_{2*}\bigl(Lp_1^*\mathcal K\otimes_E^L\mathcal K_\psi\bigr)[r],
 \end{aligned}
\]
where the second and third lines give respectively the \(!\)- and
\(*\)-transform.  There are canonical natural isomorphisms
\[
 \operatorname{For}_{V^\vee}
   \FT^{V_0/S_0}_{\psi,!}
 \simeq
 \FT^{V/S}_{\psi,!}\operatorname{For}_V,
 \qquad
 \operatorname{For}_{V^\vee}
   \FT^{V_0/S_0}_{\psi,*}
 \simeq
 \FT^{V/S}_{\psi,*}\operatorname{For}_V.
\]
They are compatible with the forget-supports transformation, and
compact-support base change transports the strict Weil datum to geometric
stalks and compactly supported cohomology.  Applying these functors at a
common finite coefficient stage and then to the idempotent in each Karoubi
triple defines
\[
 \FT^{V_0/S_0}_{\psi,!},\ \FT^{V_0/S_0}_{\psi,*}:
 \mathcal D^b_{W,c}(V_0,\Qlb)\longrightarrow
 \mathcal D^b_{W,c}(V_0^\vee,\Qlb),
\]
with the same forgetful and base-change compatibilities.
\end{theorem}

\section{Finiteness, stalks, and base change}

\begin{theorem}[Bounded constructibility]
For every \(K\in\Dbc(V,E)\), the integrand
\(\mathcal A_\psi(K)\) belongs to
\(\Dbc(V\times_SV^\vee,E)\), and compactly supported direct image gives
\[
 \FT^{V/S}_{\psi,!}:\Dbc(V,E)\longrightarrow\Dbc(V^\vee,E).
\]
The ordinary direct-image finiteness theorem for the separated finite-type
morphism \(p_2\) gives
\[
 \FT^{V/S}_{\psi,*}:\Dbc(V,E)\longrightarrow\Dbc(V^\vee,E).
\]
These restrictions are triangulated functors.  The assertions remain valid
after replacing \(E\) by a finite extension and after replacing \(V\) by its
pullback along a morphism between schemes in the geometric scope fixed
above.
\end{theorem}

\begin{definition}[Geometric fibers and compactly supported cohomology]
For a geometric point \(\bar s:\Spec\Omega\to S\), define
\[
 V_{\bar s}=V\times_S\Spec\Omega
\]
and let \(i_{\bar s}:V_{\bar s}\to V\) be the projection.  For
\(K\in\Db(V,E)\), its derived restriction to the geometric fiber is
\[
 K_{\bar s}=Li_{\bar s}^*K\quad\text{in }\Db(V_{\bar s},E).
\]
If \(Z\) is separated and of finite type over the algebraically closed
field \(\Omega\), with structural morphism
\(a_Z:Z\to\Spec\Omega\), define
\[
 R\Gamma_c(Z,L)=Ra_{Z!}L\quad\text{in }D(E),
 \qquad
 H_c^i(Z,L)=H^i(R\Gamma_c(Z,L)).
\]
For a geometric point \(i_{\bar y}:\Spec\Omega\to Y\), the derived stalk
notation \(M_{\bar y}\) means \(Li_{\bar y}^*M\) in \(D(E)\).
\end{definition}

\begin{theorem}[Geometric stalk formula]
For \(K\in\Dbc(V,E)\), let \(\bar\xi\to V^\vee\) be a geometric point with
image \(\bar s\to S\).  If \(K_{\bar s}\) denotes the pullback of \(K\) to
the geometric fiber \(V_{\bar s}\), then
\[
 \bigl(\FT^{V/S}_{\psi,!}K\bigr)_{\bar\xi}
 \simeq
 R\Gamma_c\!\left(
 V_{\bar s},
 K_{\bar s}\otimes_E^L
 \langle- ,\bar\xi\rangle^*\mathcal L_\psi
 \right)[r].
\]
The isomorphism is functorial in \(K\).  It is Frobenius-equivariant under
the following precise additional hypotheses: \(S,V,K\) are the geometric
pullbacks of \(S_0,V_0\) and a strict Weil complex \(\mathcal K\), and
\(\xi\in V_0^\vee(\mathbf F_{q^n})\) has image
\(s\in S_0(\mathbf F_{q^n})\).  Then the left side carries
\(F_{\xi,n}\), while the right side carries the operator obtained from the
\(n\)-fold strict datum on the fiber and the integrand followed by
\(R\Gamma_c\), and the displayed isomorphism intertwines these operators.
No Frobenius action is asserted for an arbitrary geometric point without
such descended rational data.
\end{theorem}

\begin{theorem}[Base change and change of coefficients]
Let \(g:S'\to S\), let \(V'=V\times_SS'\), and let
\(g_V:V'\to V\) and \(g_{V^\vee}:(V')^\vee\to V^\vee\) be the induced maps.
For every \(K\in\Dbc(V,E)\), compact-support base change gives a natural
isomorphism in \(\Dbc((V')^\vee,E)\)
\[
 Lg_{V^\vee}^*\FT^{V/S}_{\psi,!}(K)
 \simeq
 \FT^{V'/S'}_{\psi,!}\bigl(Lg_V^*K\bigr).
\]
These isomorphisms satisfy the identity and composition cocycles.  The
corresponding statement for \(\FT_{\psi,*}\) holds whenever the ordinary
base-change morphism for \(Rp_{2*}\) is an isomorphism on
\(\mathcal A_\psi(K)\); no unrestricted ordinary-direct-image base-change
claim is implicit.

If \(E'/E\) is a finite extension containing the values of \(\psi\), then
\[
 E'\otimes_E^L\FT^{V/S}_{\psi,!}(K)
 \simeq
 \FT^{V/S}_{\psi,!}\bigl(E'\otimes_E^LK\bigr),
\]
and likewise for \(\FT_{\psi,*}\) on its finiteness domain.  These
comparisons are transitive in \(E'/E\) and define the corresponding
\(\Qlb\)-coefficient transforms by passage through finite extensions of
\(\Ql\) over which the objects and morphisms are defined.
\end{theorem}

\begin{corollary}[Trace-function normalization]
Let \(S_0=\Spec k_0\), let \(V_0\) be an \(r\)-dimensional vector space
over \(k_0\), and write \(V=V_0\times_{k_0}\bar k\).  For
\(K_0\in\Dbc(V_0,E)\), let \(\mathcal K\) be its strict Weil pullback and
let \(K=\operatorname{For}_V(\mathcal K)\).  Compact-support base change
identifies the geometric pullback of \(\FT^{V_0/S_0}_{\psi,!}K_0\) with
the strict Weil Fourier transform of \(\mathcal K\).  For
\(v\in V_0(\mathbf F_{q^n})\), define
\[
 t_{K_0}^{(n)}(v)=\sum_i(-1)^i
 \Tr\!\left(F_{v,n}\mid\mathcal H^i(K)_{\bar v}\right).
\]
Then, for \(\xi\in V_0^\vee(\mathbf F_{q^n})\),
\[
 t_{\FT^{V_0/S_0}_{\psi,!}K_0}^{(n)}(\xi)
 =(-1)^r\sum_{v\in V_0(\mathbf F_{q^n})}
 t_{K_0}^{(n)}(v)\,
 \psi\!\left(
 \Tr_{\mathbf F_{q^n}/\Fq}\langle v,\xi\rangle
 \right).
\]
\end{corollary}

\section{Fourier inversion}

The evaluation map identifies \(V^{\vee\vee}\) canonically with \(V\).  Let
\(a_V:V\to V\) be multiplication by \(-1\).

\begin{definition}[Relative additive-character integral]
Let
\[
 \pi:V\times_SV^\vee\longrightarrow V,
 \qquad h(v,\xi)=\langle v,\xi\rangle,
\]
and let \(0:S\hookrightarrow V\) be the zero section.  The relative
additive-character integral is
\[
 \mathcal O_\psi=R\pi_!h^*\mathcal L_{\psi,S}
 \quad\text{in }\Dbc(V,E).
\]
\end{definition}

\begin{lemma}[Orthogonality on a line]
For the rank-one bundle \(\Aone_S\), restriction of \(\mathcal O_\psi\)
to \(\mathbf G_{m,S}\) is zero, while its pullback to the zero section is
\(E_S(-1)[-2]\).  The adjunction morphism from extension by zero therefore
gives a natural isomorphism
\[
 \mathcal O_\psi\xrightarrow{\sim}R0_!E_S(-1)[-2].
\]
The isomorphism is stable under arbitrary base change on \(S\), finite
extension of coefficients, and strict Weil descent.
\end{lemma}

\begin{proposition}[Relative orthogonality in rank \(r\)]
For a vector bundle \(V/S\) of rank \(r\), there is a canonical
isomorphism
\[
 R\pi_!h^*\mathcal L_{\psi,S}
 \xrightarrow{\sim}R0_!E_S(-r)[-2r].
\]
It is compatible with direct sums of vector bundles, changes of local
trivialization, arbitrary base change on \(S\), and strict Weil structures.
Equivalently, every nonzero fiberwise linear functional has zero
compactly supported Artin--Schreier cohomology, while the zero functional
has compactly supported cohomology \(E(-r)[-2r]\).
\end{proposition}

\begin{definition}[Double-transform kernels]
On
\(Q=V\times_SV^\vee\times_SV\), let
\(q_{12},q_{23},q_{13}\) be the evident projections.  Put
\[
 \mathcal K_\chi^{12}
 =\langle\ ,\ \rangle_{V,V^\vee}^*\mathcal L_\chi
 \quad\text{on }V\times_SV^\vee,
 \qquad
 \mathcal K_\chi^{23}
 =\langle\ ,\ \rangle_{V^\vee,V}^*\mathcal L_\chi
 \quad\text{on }V^\vee\times_SV.
\]
Thus the two factors below live on the respective targets of \(q_{12}\)
and \(q_{23}\).  Define
\[
 \begin{aligned}
 \mathcal C_+&=Rq_{13!}\bigl(
 q_{12}^*\mathcal K_\psi^{12}\otimes_E^L
 q_{23}^*\mathcal K_\psi^{23}\bigr)[2r],\\
 \mathcal C_-&=Rq_{13!}\bigl(
 q_{12}^*\mathcal K_\psi^{12}\otimes_E^L
 q_{23}^*\mathcal K_{\psi^{-1}}^{23}\bigr)[2r].
 \end{aligned}
\]
Let
\[
 \Delta(v)=(v,v),\qquad \Delta_-(v)=(v,-v)
\]
be the diagonal and negative diagonal in \(V\times_SV\).
\end{definition}

\begin{lemma}[Phase-addition comparison]
The multiplicative structure of \(\mathcal L_\psi\) gives canonical
isomorphisms on \(Q\)
\[
 q_{12}^*\mathcal K_\psi^{12}\otimes q_{23}^*\mathcal K_\psi^{23}
 \simeq
 \mathcal L_\psi(\langle v+w,\xi\rangle),
\]
and
\[
 q_{12}^*\mathcal K_\psi^{12}\otimes
 q_{23}^*\mathcal K_{\psi^{-1}}^{23}
 \simeq
 \mathcal L_\psi(\langle v-w,\xi\rangle).
\]
These isomorphisms satisfy the associativity and symmetry constraints
inherited from \(c_\psi\).
\end{lemma}

\begin{theorem}[Diagonal and antidiagonal kernels]
There are canonical isomorphisms in \(\Dbc(V\times_SV,E)\)
\[
 \mathcal C_+\xrightarrow{\sim}R\Delta_{-!}E_V(-r),
 \qquad
 \mathcal C_-\xrightarrow{\sim}R\Delta_!E_V(-r).
\]
These isomorphisms are compatible with the phase-addition comparison for
the fiberwise linear maps \((v,w)\mapsto v+w\) and \((v,w)\mapsto v-w\),
and with relative orthogonality in rank \(r\).  The cohomological normalization on
the right is \([-2r+2r]=[0]\).
Both isomorphisms commute with base change, coefficient extension, and
strict Weil descent.
\end{theorem}

\begin{proposition}[Composition of integral transforms]
For \(K\in\Dbc(V,E)\), there are functorial isomorphisms
\[
 \FT^{V^\vee/S}_{\psi,!}\FT^{V/S}_{\psi,!}(K)
 \simeq
 R\operatorname{pr}_{2!}
 \bigl(L\operatorname{pr}_1^*K\otimes_E^L\mathcal C_+\bigr),
\]
and
\[
 \FT^{V^\vee/S}_{\psi^{-1},!}\FT^{V/S}_{\psi,!}(K)
 \simeq
 R\operatorname{pr}_{2!}
 \bigl(L\operatorname{pr}_1^*K\otimes_E^L\mathcal C_-\bigr),
\]
where \(\operatorname{pr}_1,\operatorname{pr}_2:V\times_SV\to V\).
These isomorphisms are compatible with proper-support Fubini,
compact-support composition, the projection formula, and exchange of
tensor factors.  They are natural in \(K\) and compatible with distinguished
triangles.
\end{proposition}

\begin{theorem}[Fourier inversion]
For every \(K\in\Dbc(V,E)\), there are functorial isomorphisms
\[
 \FT^{V^\vee/S}_{\psi^{-1},!}
 \bigl(\FT^{V/S}_{\psi,!}K\bigr)
 \simeq K(-r),
\]
and
\[
 \FT^{V^\vee/S}_{\psi,!}
 \bigl(\FT^{V/S}_{\psi,!}K\bigr)
 \simeq a_V^*K(-r).
\]
There is no residual cohomological shift.  Consequently
\(\FT^{V/S}_{\psi,!}\) is an equivalence with quasi-inverse
\[
 K'\longmapsto
 \FT^{V^\vee/S}_{\psi^{-1},!}(K')(r).
\]
The inversion isomorphisms are compatible with base change, finite scalar
extension, and Weil structures; their unit and counit satisfy the triangle
identities.
\end{theorem}

\begin{corollary}[Equivalence on strict Weil categories]
For a vector bundle \(V_0/S_0\) over \(\Fq\), the strict-Weil transform
\[
 \FT^{V_0/S_0}_{\psi,!}:\mathcal D^b_{W,c}(V_0,E)
 \longrightarrow\mathcal D^b_{W,c}(V_0^\vee,E)
\]
is an exact triangulated equivalence.  Its quasi-inverse is
\(\FT^{V_0^\vee/S_0}_{\psi^{-1},!}(r)\).  The same assertion holds over
every finite coefficient extension and on the common-stage Karoubi category
with \(\Qlb\)-coefficients.
\end{corollary}

\section{Verdier duality and the clean condition}

Retain the Verdier-admissibility data of
Definition~\ref{def:verdier-admissible}.  Thus
\[
 \omega_X=Ra_X^!E_k\quad\text{in }\Dbc(X,E),
\]
and the contravariant exact triangulated functor used below is
\[
 D_X(K)=R\mathcal Hom_X(K,\omega_X):
 \Dbc(X,E)^{\mathrm{op}}\longrightarrow\Dbc(X,E).
\]

\begin{definition}[Fourier-clean object]
An object \(K\in\Dbc(V,E)\) is \emph{\(\psi\)-Fourier-clean} if the canonical
map
\[
 \alpha_K:\FT_{\psi,!}(K)\longrightarrow\FT_{\psi,*}(K)
\]
is an isomorphism.  To express this condition through the localization
triangle, choose a compactification
\[
 U=V\times_SV^\vee\xrightarrow{\,j\,}\overline P
 \xrightarrow{\,\bar p_2\,}V^\vee
\]
with \(j\) open and \(\bar p_2\) proper, let
\(i:Z=\overline P\setminus U\hookrightarrow\overline P\) be the closed
complement, and define the canonical boundary object
\[
 \partial_{\overline P}\mathcal A_\psi(K)
 =Ri_*Li^*Rj_*\mathcal A_\psi(K)\quad\text{in }\Dbc(\overline P,E).
\]
The localization theorem supplies the functorial distinguished triangle
\[
 Rj_!\mathcal A_\psi(K)\longrightarrow
 Rj_*\mathcal A_\psi(K)\longrightarrow
 \partial_{\overline P}\mathcal A_\psi(K)\longrightarrow.
\]
After applying \(R\bar p_{2*}(-)[r]\), this is the canonical distinguished
triangle whose first arrow is \(\alpha_K\).  Consequently \(K\) is clean
if and only if
\[
 R\bar p_{2*}\partial_{\overline P}\mathcal A_\psi(K)=0.
\]
This criterion is independent of the compactification because cleanliness
was defined by the canonical morphism \(\alpha_K\).

The object \(K\) is \emph{universally \(\psi\)-Fourier-clean over \(S\)}
if, for every morphism \(g:T\to S\) with \(T\) noetherian of finite Krull
dimension and \(\ell\) invertible on \(T\), the pullback
\(K_T=Lg_V^*K\) on \(V_T=V\times_ST\) has
\[
 \alpha_{K_T}:\FT^{V_T/T}_{\psi,!}(K_T)
 \xrightarrow{\ \sim\ }\FT^{V_T/T}_{\psi,*}(K_T).
\]
Equivalently, the displayed boundary direct image vanishes for the base
change to \(T\) of any compactification as above.  For an arbitrary
integral transform, cleanliness is an additional condition.  For the
Artin--Schreier Fourier kernel used in this note, it is automatic by the
universal boundary-acyclicity theorem below.
\end{definition}

\begin{definition}[Universal local acyclicity]
Let \(f:X\to Y\) be a morphism and let \(L\in\Dbc(X,E)\).  For a geometric
point \(\bar x\to X\), with image \(\bar y\to Y\), write
\[
 X_{(\bar x)}=\Spec\mathcal O^{\mathrm{sh}}_{X,\bar x},
 \qquad
 Y_{(\bar y)}=\Spec\mathcal O^{\mathrm{sh}}_{Y,\bar y}.
\]
The pair \((f,L)\) is \emph{locally acyclic} if, for every such \(\bar x\)
and every geometric point \(\bar t\to Y_{(\bar y)}\) specializing to
\(\bar y\), the canonical specialization morphism
\[
 L_{\bar x}\longrightarrow
 R\Gamma\!\left(
 X_{(\bar x)}\times_{Y_{(\bar y)}}\bar t,
 L|_{X_{(\bar x)}\times_{Y_{(\bar y)}}\bar t}
 \right)
\]
is an isomorphism in \(D^b_{\mathrm{fd}}(E)\).  It is \emph{universally
locally acyclic}, abbreviated ULA, if the same condition holds after every
base change \(Y'\to Y\).  Equivalently, all corresponding vanishing-cycle
complexes are zero after every base change.
\end{definition}

\begin{theorem}[Universal local acyclicity at the Fourier boundary]
Assume \(S=\Spec k\), let \(V=\Aone_k\), let \(V^\vee=\Aone_k\) be its
dual line, and let the pairing be \((x,x')\mapsto xx'\).  Put
\[
 J:\Aone_k\times_k\Aone_k
   \hookrightarrow\mathbf P^1_k\times_k\Aone_k,
 \qquad
 \overline p_1:\mathbf P^1_k\times_k\Aone_k
   \longrightarrow\mathbf P^1_k.
\]
Then
\[
 \left(\overline p_1,RJ_!\mathcal K_\psi\right)
\]
is universally locally acyclic.  Equivalently, its vanishing cycles at
\(\infty\in\mathbf P^1_k\) are zero after every base change on
\(\mathbf P^1_k\).  The same assertion holds with \(\psi^{-1}\) and after
finite extension of the coefficient field.  In the arithmetic setting the
ULA isomorphisms commute with the strict Weil structures.
\end{theorem}

\begin{theorem}[Laumon's universal \(!\)-to-\(*\) comparison]
For a vector bundle \(V/S\) of rank \(r\), put
\[
 \overline V=
 \operatorname{Proj}_S
 \operatorname{Sym}_{\mathcal O_S}
 \bigl(\mathcal O_S\oplus\mathcal V^\vee\bigr),
 \qquad
 J:V\times_SV^\vee\hookrightarrow\overline V\times_SV^\vee,
\]
where the degree-one generator from \(\mathcal O_S\) defines the displayed
open immersion, and let
\(\overline p_2:\overline V\times_SV^\vee\to V^\vee\) be the second
projection.  For every \(K\in\Dbc(V,E)\), every morphism \(g:T\to S\), and
\(K_T=Lg_V^*K\), one has
\[
 R\overline p_{2,T*}
 \partial_{\overline V_T\times_TV_T^\vee}
 \mathcal A_{\psi,T}(K_T)=0.
\]
These vanishings are compatible with identities and composition of base
changes.  Hence every bounded constructible complex is universally
\(\psi\)-Fourier-clean, and the canonical morphism
\[
 \alpha_K:\FT_{\psi,!}(K)\xrightarrow{\ \sim\ }\FT_{\psi,*}(K)
\]
is a functorial isomorphism.  The same statements hold for \(\psi^{-1}\)
and for strict Weil objects.
\end{theorem}

\begin{theorem}[Verdier duality for Fourier transform]
For every \(K\in\Dbc(V,E)\) such that the ordinary transform of \(D_VK\)
lies in \(\Dbc(V^\vee,E)\), there is a functorial isomorphism
\[
 D_{V^\vee}\bigl(\FT^{V/S}_{\psi,!}K\bigr)
 \simeq
 \FT^{V/S}_{\psi^{-1},*}\bigl(D_VK\bigr)(r).
\]
This is the general duality statement and requires no clean-support
hypothesis.  If \(D_VK\) is \(\psi^{-1}\)-Fourier-clean, it specializes to
\[
 D_{V^\vee}\bigl(\FT^{V/S}_{\psi,!}K\bigr)
 \simeq
 \FT^{V/S}_{\psi^{-1},!}\bigl(D_VK\bigr)(r).
\]
The symbol \((r)\) is a Tate twist.  Neither formula contains an additional
cohomological shift.
\end{theorem}

\begin{definition}[Stalk and costalk at a geometric point]
Let \(i_{\bar x}:\bar x=\Spec\Omega\to X\) be a geometric point with image
\(x\).  Let
\[
 X_{(\bar x)}=\Spec\mathcal O^{\mathrm{sh}}_{X,\bar x},
 \qquad
 u_{\bar x}:X_{(\bar x)}\longrightarrow X,
 \qquad
 s_{\bar x}:\bar x\hookrightarrow X_{(\bar x)}
\]
be the strict localization, its canonical morphism, and its closed geometric
point.  For \(K\in\Dbc(X,E)\), define the derived stalk and costalk by
\[
 i_{\bar x}^*K=Li_{\bar x}^*K,
 \qquad
 i_{\bar x}^!K=Rs_{\bar x}^!Lu_{\bar x}^*K.
\]
Here \(Rs_{\bar x}^!\) is the right adjoint of
\(Rs_{\bar x*}\) for the closed immersion; equivalently, the second object
is the local cohomology of \(Lu_{\bar x}^*K\) supported at the closed point,
viewed as a complex of \(E\)-vector spaces.  Both objects lie in
\(D^b_{\mathrm{fd}}(E)\), the full subcategory of \(D(E)\) whose cohomology
is finite-dimensional and vanishes outside a finite interval.
\end{definition}

\section{The middle perverse heart}

Let \(X\) be a finite-type scheme over a field.  For a geometric point
\(i_{\bar x}:\bar x\to X\), put
\(d(x)=\dim\overline{\{x\}}\), where \(x\) is its image.  The middle
perverse \(t\)-structure on \(\Dbc(X,E)\) is defined by the following two
full subcategories, using the stalk and costalk just defined:
\[
 \begin{aligned}
 {}^p\Db^{\le0}(X,E)
 &=\{K:\mathcal H^i(i_{\bar x}^*K)=0
       \text{ for every }x\text{ and }i>-d(x)\},\\
 {}^p\Db^{\ge0}(X,E)
 &=\{K:\mathcal H^i(i_{\bar x}^!K)=0
       \text{ for every }x\text{ and }i<-d(x)\}.
 \end{aligned}
\]
Its heart
\[
 \Perv(X,E)={}^p\Db^{\le0}(X,E)\cap{}^p\Db^{\ge0}(X,E)
\]
is the abelian category of middle perverse \(E\)-sheaves.

\begin{theorem}[Bounded middle-perverse \(t\)-structure and truncations]
The pair
\[
 \left({}^p\Db^{\le0}(X,E),{}^p\Db^{\ge0}(X,E)\right)
\]
is a bounded \(t\)-structure on the ordinary triangulated category
\(\Dbc(X,E)\).  For every \(n\in\mathbf Z\) it has truncation functors
\[
 {}^p\tau_{\le n},{}^p\tau_{\ge n}:\Dbc(X,E)\longrightarrow\Dbc(X,E)
\]
and a functorial distinguished triangle
\[
 {}^p\tau_{\le n}K\longrightarrow K\longrightarrow
 {}^p\tau_{\ge n+1}K\longrightarrow{}^p\tau_{\le n}K[1].
\]
The perverse cohomology objects are
\[
 {}^p\mathcal H^n(K)={}^p\mathcal H^0(K[n])\in\Perv(X,E),
 \qquad
 {}^p\mathcal H^0={}^p\tau_{\le0}\,{}^p\tau_{\ge0}.
\]
The heart \(\Perv(X,E)\) is artinian and noetherian.  Verdier duality
exchanges \({}^p\Db^{\le0}\) and \({}^p\Db^{\ge0}\), commutes with
perverse cohomology in the form
\[
 D_X\,{}^p\mathcal H^n(K)\simeq{}^p\mathcal H^{-n}(D_XK),
\]
and restricts to an exact contravariant autoequivalence of the perverse
heart.
\end{theorem}

\begin{theorem}[Support and cosupport criteria]
Adopt the convention \(\dim\varnothing=-\infty\).  For
\(K\in\Dbc(X,E)\), the stalk definition of the middle-perverse
\(t\)-structure is equivalent to
\[
 K\in{}^p\Db^{\le0}(X,E)
 \quad\Longleftrightarrow\quad
 \dim\operatorname{Supp}\mathcal H^i(K)\le-i
 \quad\text{for every }i.
\]
The costalk definition is equivalent to the Verdier-dual support condition
\[
 K\in{}^p\Db^{\ge0}(X,E)
 \quad\Longleftrightarrow\quad
 \dim\operatorname{Supp}\mathcal H^i(D_XK)\le-i
 \quad\text{for every }i.
\]
\end{theorem}

\begin{theorem}[Open--closed recollement and middle extension]
Let \(i:Z\hookrightarrow X\) be a closed immersion with complementary open
immersion \(j:U\hookrightarrow X\).  The middle-perverse \(t\)-structure on
\(X\) is the recollement of those on \(Z\) and \(U\).  More explicitly,
\(Ri_*=Ri_!\) is \(t\)-exact and fully faithful,
\(Lj^*=Rj^!\) is \(t\)-exact, and
\[
 \begin{aligned}
 Rj_!({}^p\Db^{\le0}(U,E))&\subset{}^p\Db^{\le0}(X,E),&
 Rj_*({}^p\Db^{\ge0}(U,E))&\subset{}^p\Db^{\ge0}(X,E),\\
 Li^*({}^p\Db^{\le0}(X,E))&\subset{}^p\Db^{\le0}(Z,E),&
 Ri^!({}^p\Db^{\ge0}(X,E))&\subset{}^p\Db^{\ge0}(Z,E).
 \end{aligned}
\]
For \(P\in\Perv(U,E)\), define
\[
 {}^pj_!P={}^p\mathcal H^0(Rj_!P),
 \qquad
 {}^pj_*P={}^p\mathcal H^0(Rj_*P),
 \qquad
 j_{!*}P=\operatorname{im}({}^pj_!P\to{}^pj_*P).
\]
The object \(j_{!*}P\) is the unique perverse extension of \(P\) having no
nonzero perverse subobject and no nonzero perverse quotient supported on
\(Z\).  The recollement functors and the intermediate extension commute
with finite extension of coefficients.
\end{theorem}

\begin{theorem}[Smooth pullback, lisse tensor product, and twists]
If \(f:Y\to X\) is smooth of pure relative dimension \(d\), then
\[
 Lf^*[d]:\Perv(X,E)\longrightarrow\Perv(Y,E)
\]
is exact.  If \(\mathcal L\) is lisse of finite rank on \(X\), then
\(K\mapsto K\otimes_E\mathcal L\) is perverse \(t\)-exact.  Every Tate
twist \(K\mapsto K(m)\) is also perverse \(t\)-exact.
\end{theorem}

\begin{theorem}[Artin--Gabber affine exactness]
Let \(f:X\to Y\) be an affine morphism of finite-type schemes over a field.
Then
\[
 Rf_!({}^p\Db^{\ge0}(X,E))\subset{}^p\Db^{\ge0}(Y,E),
 \qquad
 Rf_*({}^p\Db^{\le0}(X,E))\subset{}^p\Db^{\le0}(Y,E).
\]
Equivalently, \(Rf_!\) is left perverse \(t\)-exact and \(Rf_*\) is right
perverse \(t\)-exact.  For \(Y=\Spec k\) and \(P\in\Perv(X,E)\), these
inclusions give
\[
 H_c^i(X,P)=0\quad(i<0),
 \qquad
 H^i(X,P)=0\quad(i>0).
\]
\end{theorem}

\begin{proposition}[Perverse Fourier integrand]
Assume \(S=\Spec k\), where \(k\) is finite or algebraically closed, and let
\(P\in\Perv(V,E)\).  The complex
\[
 \mathcal B_\psi(P)=Lp_1^*P[r]\otimes_E\mathcal K_\psi
 \quad\text{in }\Dbc(V\times_kV^\vee,E)
\]
is perverse.  With this notation,
\[
 \FT_{\psi,!}(P)=Rp_{2!}\mathcal B_\psi(P),
 \qquad
 \FT_{\psi,*}(P)=Rp_{2*}\mathcal B_\psi(P).
\]
\end{proposition}

\begin{proposition}[Affine half-bounds for the two Fourier images]
Under the preceding hypotheses, the affine projection
\(p_2:V\times_kV^\vee\to V^\vee\) gives
\[
 \FT_{\psi,!}(P)\in{}^p\Db^{\ge0}(V^\vee,E),
 \qquad
 \FT_{\psi,*}(P)\in{}^p\Db^{\le0}(V^\vee,E).
\]
\end{proposition}

\begin{theorem}[Fourier support inequality]
For every \(P\in\Perv(V,E)\) and every \(i\in\mathbf Z\), one has
\[
 \dim\operatorname{Supp}\mathcal H^i(\FT_{\psi,!}P)\le-i.
\]
The same inequality holds for \(\FT_{\psi^{-1},!}P\).
\end{theorem}

\begin{theorem}[Fourier cosupport inequality]
For every \(P\in\Perv(V,E)\) and every \(i\in\mathbf Z\), one has
\[
 \dim\operatorname{Supp}
 \mathcal H^i\!\left(D_{V^\vee}\FT_{\psi,!}P\right)\le-i.
\]
The same inequality holds for \(\FT_{\psi^{-1},!}P\).
\end{theorem}

\begin{theorem}[Perverse exactness over a field]
Assume \(S=\Spec k\), where \(k\) is finite or algebraically closed.  The
normalized transform restricts to an exact
equivalence of abelian categories
\[
 \FT_{\psi,!}:\Perv(V,E)
 \xrightarrow{\ \sim\ }
 \Perv(V^\vee,E),
\]
with quasi-inverse \(\FT_{\psi^{-1},!}(r)\).  It preserves simple objects and
Jordan--Holder length.  When \(k=\Fq\), under extension of the ground field
to \(\mathbf F_{q^m}\), its base change is the transform formed with
\(\psi_m=\psi\circ\Tr_{\mathbf F_{q^m}/\Fq}\); after base change to
\(\bar k\), it induces the corresponding equivalence on the geometric
perverse hearts.

For a general base \(S\), this statement does not assert exactness for the
ordinary middle perverse hearts of the total spaces.  A relative version
requires a separately specified relative perverse \(t\)-structure together
with universal local acyclicity over \(S\).
\end{theorem}

\begin{theorem}[Perverse exactness for strict Weil objects]
Assume \(S_0=\Spec\Fq\), and define
\[
 \Perv_W(V_0,E)=
 \{\mathcal K\in\mathcal D^b_{W,c}(V_0,E):
   \operatorname{For}_V(\mathcal K)\in\Perv(V,E)\}.
\]
With morphisms inherited from the strict-Weil derived category, this is an
abelian category.  The strict-Weil Fourier transform is perverse
\(t\)-exact and induces an exact equivalence
\[
 \FT^{V_0/S_0}_{\psi,!}:\Perv_W(V_0,E)
 \xrightarrow{\sim}\Perv_W(V_0^\vee,E),
\]
whose quasi-inverse is
\(\FT^{V_0^\vee/S_0}_{\psi^{-1},!}(r)\).  Forgetting the Weil structure
commutes with this equivalence and with perverse kernels, images, cokernels,
and finite Jordan--Holder filtrations.  The same construction at common
finite coefficient stages defines the perverse hearts and equivalence over
\(\Qlb\).
\end{theorem}

\begin{proposition}[Punctual-to-exponential comparison]
Let \(V=\Aone_{\bar k}\), let \(i_a:\{a\}\hookrightarrow V\), and let
\(M\) be a one-dimensional \(\Qlb\)-vector space.  On the dual line, with
coordinate \(x'\), there is a functorial isomorphism
\[
 \FT_{\psi,!}(Ri_{a!}M)
 \simeq\mathcal L_{\psi(ax')}[1]\otimes M.
\]
The comparison commutes with finite extension of coefficients and with
strict Weil descent when \(a\) and \(M\) descend to a finite ground field.
\end{proposition}

\begin{proposition}[Exponential-to-punctual inverse comparison]
For \(a\in\bar k\) and a one-dimensional \(M\), there is a functorial
isomorphism
\[
 \FT_{\psi,!}\bigl(\mathcal L_{\psi(ax)}[1]\otimes M\bigr)
 \simeq Ri_{-a,!}M(-1).
\]
After application of \(\FT_{\psi^{-1},!}(1)\), the two exceptional
comparison isomorphisms agree with the unit and counit of Fourier inversion.
They commute with finite scalar extension and, whenever the data descend,
with strict Weil structures.
\end{proposition}

\begin{definition}[The three Fourier types on a line]
Let \(V=\Aone_{\bar k}\).  A simple object \(K\in\Perv(V,\Qlb)\) is
called
\begin{enumerate}
\item \emph{type 1} if \(K\simeq Ri_{a!}M\) for a closed point
      \(i_a:\{a\}\hookrightarrow V\) and a one-dimensional
      \(\Qlb\)-vector space \(M\);
\item \emph{type 2} if
      \(K\simeq\mathcal L_{\psi(ax)}[1]\otimes M\) for some
      \(a\in\bar k\) and a one-dimensional constant \(M\);
\item \emph{type 3} if it is neither type 1 nor type 2.
\end{enumerate}
A simple object of \(\Perv_W(\Aone_0,E)\) is \emph{arithmetically type 3}
if every simple constituent of its geometric pullback after scalar
extension to \(\Qlb\) is geometrically type 3.
\end{definition}

\begin{theorem}[Trichotomy and the full-support case]
Every simple object of \(\Perv(\Aone_{\bar k},\Qlb)\) belongs to exactly
one of the three types.  A simple object has full support if and only if it
has the form
\[
 j_{!*}(\mathcal F[1])
\]
for a dense open immersion \(j:U\hookrightarrow\Aone_{\bar k}\) and an
irreducible lisse sheaf \(\mathcal F\) on \(U\).  Such an object is type 2
exactly when, on some dense open, \(\mathcal F\) is isomorphic to
\(\mathcal L_{\psi(ax)}\otimes M\) for some \(a\in\bar k\) and some
one-dimensional constant \(M\); otherwise it is type 3.
\end{theorem}

\begin{theorem}[Simplicity of a Fourier image]
If \(K\in\Perv(\Aone_{\bar k},\Qlb)\) is simple, then
\(\FT_{\psi,!}(K)\) is simple.  The same assertion holds for the inverse
transform \(\FT_{\psi^{-1},!}(1)\).
\end{theorem}

\begin{lemma}[Inverse detection of a punctual Fourier image]
For every simple \(K\in\Perv(\Aone_{\bar k},\Qlb)\),
\[
 \FT_{\psi,!}(K)\text{ is of type 1}
 \quad\Longleftrightarrow\quad
 K\text{ is of type 2}.
\]
Under this equivalence the parameters and coefficient lines are related by
the two exceptional comparison isomorphisms and the Tate twist in Fourier
inversion.
\end{lemma}

\begin{lemma}[Inverse detection of an exponential Fourier image]
For every simple \(K\in\Perv(\Aone_{\bar k},\Qlb)\),
\[
 \FT_{\psi,!}(K)\text{ is of type 2}
 \quad\Longleftrightarrow\quad
 K\text{ is of type 1}.
\]
Under this equivalence the support point and the linear phase are related by
the sign in the exponential-to-punctual comparison.
\end{lemma}

\begin{corollary}[Exclusion of type 1 for a type-3 input]
If \(K\) is simple of type 3, then \(\FT_{\psi,!}(K)\) is not of type 1.
\end{corollary}

\begin{corollary}[Exclusion of type 2 for a type-3 input]
If \(K\) is simple of type 3, then \(\FT_{\psi,!}(K)\) is not of type 2.
\end{corollary}

\begin{corollary}[Type-3 preservation and full support]
If \(K\) is simple of type 3, then \(\FT_{\psi,!}(K)\) is simple of type 3
and has full support.
\end{corollary}

\begin{theorem}[Maximal lisse locus of a type-3 transform]
Let \(K=j_{!*}(\mathcal F[1])\) be simple of type 3 and put
\(K'=\FT_{\psi,!}(K)\).  There is a unique maximal dense open immersion
\[
 j':U'\hookrightarrow(\Aone_{\bar k})^\vee
\]
on which \(j'^*K'\) is lisse up to the perverse shift.  The sheaf
\[
 \mathcal F'=\mathcal H^{-1}(j'^*K')
\]
is irreducible and lisse on \(U'\), and the canonical identification is
\[
 j'^*K'\xrightarrow{\ \sim\ }\mathcal F'[1].
\]
\end{theorem}

\begin{theorem}[Middle-extension comparisons and inverse equations]
With the preceding notation, restriction to \(U'\) determines comparison
maps
\[
 c_K:K'\longrightarrow j'_{!*}(\mathcal F'[1]),
 \qquad
 d_K:j'_{!*}(\mathcal F'[1])\longrightarrow K'
\]
whose restrictions to \(U'\) are the identity of \(\mathcal F'[1]\).  They
satisfy
\[
 d_Kc_K=\operatorname{id}_{K'},
 \qquad
 c_Kd_K=\operatorname{id}_{j'_{!*}(\mathcal F'[1])}.
\]
Thus \(K'\simeq j'_{!*}(\mathcal F'[1])\), and \(K'\) has no nonzero
punctual perverse subobject, quotient, or direct summand.
\end{theorem}

\begin{theorem}[Uniqueness of the type-3 transform datum]
Suppose \(j'':U''\hookrightarrow(\Aone_{\bar k})^\vee\) is maximal among
dense opens on which \(j''^*\FT_{\psi,!}(K)\) is lisse up to shift,
\(\mathcal F''\) is irreducible and lisse on \(U''\), and an isomorphism
\[
 \FT_{\psi,!}(K)\xrightarrow{\sim}j''_{!*}(\mathcal F''[1])
\]
is specified.  Then \(U''=U'\), and there is a unique isomorphism
\(\mathcal F''\xrightarrow{\sim}\mathcal F'\) compatible with the two
restrictions of \(\FT_{\psi,!}(K)\).  This uniqueness is preserved by
finite extension of coefficients.
\end{theorem}

\begin{theorem}[Exact type exchange and type-3 transform shape]
On the geometric perverse heart of the affine line, \(\FT_{\psi,!}\)
exchanges types 1 and 2 and preserves type 3.  For every one-dimensional
\(M\), the exceptional cases are
\[
 \FT_{\psi,!}(Ri_{a!}M)
   \simeq\mathcal L_{\psi(ax')}[1]\otimes M,
 \qquad
 \FT_{\psi,!}(\mathcal L_{\psi(ax)}[1]\otimes M)
   \simeq Ri_{-a,!}M(-1).
\]
For type 3, the transform is exactly the middle extension
\[
 \FT_{\psi,!}\bigl(j_{!*}(\mathcal F[1])\bigr)
 \simeq j'_{!*}(\mathcal F'[1])
\]
on its unique maximal lisse open, with \(\mathcal F'\) irreducible and
lisse and with no punctual perverse subobject, quotient, or direct summand.
\end{theorem}

\begin{theorem}[Arithmetic type-3 descent]
Let \(K_0\) be a simple object of \(\Perv_W(\Aone_0,E)\) that is
arithmetically type 3.  Then its strict-Weil Fourier transform is
arithmetically type 3.  After geometric pullback and extension to a common
finite coefficient field, each simple constituent has the maximal-lisse-open
and middle-extension description above.  These descriptions are compatible
with further finite coefficient extension and with the common-stage
\(\Qlb\)-construction.
\end{theorem}

\section{Convolution and tensor product}

\begin{definition}[External product and compactly supported convolution]
Let \(m:V\times_SV\to V\) be addition and let
\(q_1,q_2:V\times_SV\to V\) be the projections.  For
\(K,L\in\Dbc(V,E)\), define the external product and the unnormalized
compactly supported convolution by
\[
 K\boxtimes_S^LL=Lq_1^*K\otimes_E^LLq_2^*L,
 \qquad
 K\star_!L=Rm_!(K\boxtimes_S^LL).
\]
Let \(E_S\) denote the constant sheaf with value \(E\) on \(S\).  The zero
section \(0:S\to V\) gives the convolution unit
\(\delta_0=R0_!E_S\).  The normalized convolution is
\[
 K*L=(K\star_!L)[r];
\]
its unit is \(\delta_0[-r]\).
\end{definition}

\begin{lemma}[Fourier factorization of the addition kernel]
On \(V\times_SV\times_SV^\vee\), there is a canonical isomorphism
\[
 \mathcal L_\psi(\langle v_1+v_2,\xi\rangle)
 \simeq
 \mathcal L_\psi(\langle v_1,\xi\rangle)
 \otimes_E
 \mathcal L_\psi(\langle v_2,\xi\rangle).
\]
For all \(K,L\in\Dbc(V,E)\), the associated transform isomorphism is
\[
 \FT_{\psi,!}(Rm_!(K\boxtimes_S^LL))[r]
 \xrightarrow{\sim}
 \FT_{\psi,!}(K)\otimes_E^L\FT_{\psi,!}(L)
\]
Both isomorphisms are compatible with compact-support composition, base
change for the addition square, and the projection formula.  They are also
compatible with the associativity, symmetry, and unit constraints.
\end{lemma}

\begin{theorem}[Convolution theorem]
For \(K,L\in\Dbc(V,E)\), there are natural isomorphisms
\[
 \FT_{\psi,!}(K\star_!L)
 \simeq
 \bigl(\FT_{\psi,!}K\otimes_E^L\FT_{\psi,!}L\bigr)[-r],
\]
or, equivalently,
\[
 \FT_{\psi,!}(K*L)
 \simeq
 \FT_{\psi,!}K\otimes_E^L\FT_{\psi,!}L.
\]
Conversely, if \(\star_!^\vee\) denotes unnormalized convolution on the dual
bundle, then
\[
 \FT_{\psi,!}(K\otimes_E^LL)
 \simeq
 \bigl(\FT_{\psi,!}K\star_!^\vee
       \FT_{\psi,!}L\bigr)[r](r).
\]
These isomorphisms commute with associativity, symmetry, unit constraints,
base change, and finite extension of coefficients.  In particular, if
\(\pi:V\to S\) and \(\pi^\vee:V^\vee\to S\) are the bundle maps, then
\[
 \FT_{\psi,!}(\delta_0)\simeq L(\pi^\vee)^*E_S[r],
 \qquad
 \FT_{\psi,!}\bigl(L\pi^*M[r]\bigr)
 \simeq R0^\vee_!M(-r)
\]
for \(M\in\Dbc(S,E)\), where \(0^\vee:S\to V^\vee\) is the dual zero
section.
\end{theorem}

\section{Translation, modulation, and linear maps}

For a section \(a:S\to V\), let \(\tau_a(v)=v+a\) and
\(\ell_a(\xi)=\langle a,\xi\rangle\).  For a dual section
\(\xi_0:S\to V^\vee\), let
\(\ell_{\xi_0}(v)=\langle v,\xi_0\rangle\) and
\(\tau_{\xi_0}(\xi)=\xi+\xi_0\).

\begin{theorem}[Translation and modulation]
For every \(K\in\Dbc(V,E)\), there are natural isomorphisms in
\(\Dbc(V^\vee,E)\)
\[
 \FT_{\psi,!}(R\tau_{a,!}K)
 \simeq
 \FT_{\psi,!}(K)\otimes_E^L\ell_a^*\mathcal L_\psi,
\]
and
\[
 \FT_{\psi,!}\bigl(K\otimes_E^L
 \ell_{\xi_0}^*\mathcal L_\psi\bigr)
 \simeq
 L\tau_{\xi_0}^*\FT_{\psi,!}(K).
\]
They are compatible with addition of sections.
\end{theorem}

\begin{theorem}[Linear maps]
Let \(f:V_1\to V_2\) be an \(S\)-linear map between vector bundles of ranks
\(r_1\) and \(r_2\), and let \(f^\vee:V_2^\vee\to V_1^\vee\) be its
transpose.  For every \(K\in\Dbc(V_1,E)\), bounded constructibility of
\(Rf_!K\) and of both Fourier transforms gives an isomorphism in
\(\Dbc(V_2^\vee,E)\)
\[
 \FT^{V_2/S}_{\psi,!}(Rf_!K)
 \simeq
 L(f^\vee)^*\FT^{V_1/S}_{\psi,!}(K)[r_2-r_1].
\]
The isomorphism is functorial under composition and compatible with base
change and coefficient extension.
\end{theorem}

\section{References}

The normalizations and structural results recorded here are those of
Laumon's theory of the Fourier--Deligne transform, together with the
constructibility, base-change, and duality theorems of the etale and
pro-etale six-functor formalisms.  Standard references are SGA~4,
SGA~4\(\frac12\), SGA~5, Deligne's \emph{Weil I} and \emph{Weil II},
Laumon's \emph{Transformation de Fourier, constantes d'equations
fonctionnelles et conjecture de Weil}, and the Stacks Project chapters on
etale cohomology, derived categories, and compactly supported cohomology.

\end{document}
